% 第8章 8.1.2：线性插值样条与帽形基。结点 t_i=i/6（i=0,...,6），h=1/6。
% (a) g(x)=0.5+0.35 sin(2 pi x)（灰）与其线性插值 S_h g（蓝折线，圆点为结点值）；
% (b) 帽形函数 phi_0,...,phi_6（细灰），phi_3 加粗；把 t_3 处结点值由 g(1/2)=0.5 改为 0.8 后（黄虚线），
%     插值只在 [t_2,t_4]=[1/3,2/3] 上改变。曲线由解析式计算。
\begin{tikzpicture}[font=\figurefont\footnotesize]
\pgfplotsset{every axis/.style={width=52mm,height=36mm,axis lines=left,
  x axis line style={draw=FigureStroke,line width=.6pt},y axis line style={draw=FigureStroke,line width=.6pt},
  tick style={draw=FigureStroke},tick label style={font=\figurefont\scriptsize,text=FigureInk},
  xmin=0,xmax=1.04,xtick={0,0.3333,0.6667,1},xticklabels={$0$,$1/3$,$2/3$,$1$},
  title style={font=\figurefont\footnotesize},xlabel={$x$},label style={font=\figurefont\scriptsize},
  xlabel style={at={(axis description cs:1,0)},anchor=north west,yshift=1mm}}}
\begin{axis}[name=A,ymin=0,ymax=1,ytick={0,0.5,1},title={(a) 结点值与线性插值}]
\addplot[draw=FigureMuted,line width=.8pt,domain=0:1,samples=150] {0.5+0.35*sin(360*x)};
\addplot[draw=FigureBlue,line width=1.1pt,mark=*,mark size=1.3pt,mark options={fill=FigureBlue},samples at={0,0.16667,0.33333,0.5,0.66667,0.83333,1}] {0.5+0.35*sin(360*x)};
\end{axis}
\begin{axis}[at={(A.east)},anchor=west,xshift=12mm,ymin=0,ymax=1.25,ytick={0,0.5,1},title={(b) 帽形基与局部修改}]
\foreach \i in {0,...,6}{
  \addplot[draw=FigureMuted!55,line width=.5pt,domain=0:1,samples=61] {max(0,1-abs(6*x-\i))};}
\addplot[draw=FigureBlue,line width=1.1pt,domain=0:1,samples=61] {max(0,1-abs(6*x-3))};
\addplot[draw=FigureInk,line width=.8pt,samples at={0,0.16667,0.33333,0.5,0.66667,0.83333,1}] {0.5+0.35*sin(360*x)};
\addplot[draw=FigureYellow,line width=1.1pt,dashed] coordinates {(0.33333,0.80311) (0.5,0.8) (0.66667,0.19689)};
\node[anchor=south,font=\figurefont\scriptsize,text=FigureBlue] at (axis cs:0.5,1.01) {$\varphi_3$};
\end{axis}
\end{tikzpicture}
